\documentclass[10pt,a4paper]{amsart}
\usepackage[margin=19mm]{geometry}
\usepackage{amsmath,amssymb,mathtools}
\usepackage[scr=rsfs,cal=euler]{mathalfa}
\usepackage{xfrac}
\usepackage[unicode,hidelinks,bookmarks=false]{hyperref}
\newcommand{\ie}{i.e.\ }
\newcommand{\cf}{cf.\ }
\newcommand{\resp}{respectively\ }
\newcommand{\Subsection}{\subsection}
\providecommand{\nobreakdash}{\nobreak\hbox{-}}
\setlength{\emergencystretch}{2em}
\multlinegap0pt
\usepackage{bm}
\usepackage{bbm}
\usepackage{dsfont}
\usepackage{xspace}
\usepackage{cancel}
\usepackage{mathtools}
\usepackage{cleveref}
\usepackage{amsthm}
\makeatletter
\@ifundefined{theorem}{\newtheorem{theorem}{Theorem}}{}
\makeatother
\newcommand{\DMPN}{\mathrm{DMPN}}
\newcommand{\MPN}{\mathrm{MPN}}
\newcommand{\cgeq}{\geq_{\cone}}
\newcommand{\cone}{\mathcal{K}}
\newcommand{\cradius}{\rho_{\cone}}
\newcommand{\family}{\mathcal{F}}
\newcommand{\interior}{\mathrm{Int}}
\title[Nonlinear Joint Spectral Radius]{Nonlinear Joint Spectral Radius}
\author{Piero Deidda, Nicola Guglielmi, Francesco Tudisco}
\date{Source: arXiv:2507.11314v1 (CC BY 4.0)}
\begin{document}
\maketitle

\noindent\emph{Source and licence.} Nonlinear Joint Spectral Radius. Version arXiv:2507.11314v1 (CC BY 4.0), distributed under the Creative Commons Attribution 4.0 International licence, as recorded in the arXiv licence metadata for this exact version and shown on the abstract page https://arxiv.org/abs/2507.11314v1. Full rights notice: licenses/arxiv-2507.11314-CC-BY-4.0.txt.\par\smallskip

\noindent\textit{Editorial scope.} Counted unit: the Theorem whose source label is \texttt{Thm\_duality\_for\_the\_cone\_spectral\_radius} in the article (source0.tex lines 1170--1183, source label \texttt{Thm\_duality\_for\_the\_cone\_spectral\_radius}), delivered together with the article's own complete original proof of it (source0.tex lines 1186--1245, one \texttt{proof} environment of 60 lines). No mathematics is written, repaired or re-derived: statement and proof are the article's own text line for line. Required context retained uncounted: Cone\_joint\_spectral\_radius (lines 356--358); Theorem\_asymptotic\_stability\_via\_JSR (lines 415--417). Every definition, display and label of those lines is retained unchanged. Omitted from this package and not redistributed: 1--355, 359--414, 418--1169, 1184--1185, 1246--2819. Nothing inside a retained excerpt is omitted or altered: every retained body is delivered exactly as the article's lines are written, so no omission receipt of any kind accompanies this package. Citations: the retained excerpts carry no citation command at all, so no bibliography entry of the article is transcribed and no reference is redistributed. Reference adaptation: none was needed and none was made; every reference inside the delivered unit resolves inside it, so no unresolved reference or citation is produced. Printed slips kept literal: no printed word, symbol, hypothesis or number of the counted statement or of its proof was altered. Layout and class: the article's own class is \texttt{article} with packages including babel, authblk, color, amsmath, amssymb, mathdots, accents, amsthm, geometry, ifthen, caption, subcaption, mathtools, hyperref; the wrapper is the lane's pinned \texttt{amsart} editorial wrapper at 10pt on A4, which restates the theorem-like environments the retained text needs and adds the article's own macro definitions that the retained text needs (\texttt{\textbackslash DMPN}, \texttt{\textbackslash MPN}, \texttt{\textbackslash cgeq}, \texttt{\textbackslash cone}, \texttt{\textbackslash cradius}, \texttt{\textbackslash family}, \texttt{\textbackslash interior}), with extra wrapper packages (bm, bbm, dsfont, xspace, cancel, mathtools, cleveref). The delivered numbering is the unit's own. Assets: no retained span contains a figure environment, a graphics-inclusion command, a PDF-page inclusion, a table or a TikZ picture; no image file is copied into this package, the package declares no asset dependency and no image file is redistributed with it. Licence: version arXiv:2507.11314v1 of this article is distributed under the Creative Commons Attribution 4.0 International licence, as recorded in the arXiv licence metadata for this exact version and shown on the abstract page https://arxiv.org/abs/2507.11314v1. A full rights notice is delivered at licenses/arxiv-2507.11314-CC-BY-4.0.txt. Characters: every retained excerpt is pure ASCII, so the delivered engine, XeLaTeX with its default Latin Modern text font, reports no missing character.
\begin{equation}\label{Cone_joint_spectral_radius}
\cradius(\family)=\limsup_{k}\sup_{f\in\Sigma_k(\family)}\big(\|f\|_{\cone}\big)^{\frac{1}{k}}
\end{equation}

\begin{theorem}\label{Theorem_asymptotic_stability_via_JSR}
    Let $\family=\{f_i\}_{i\in I}$ be a bounded family of subhomogeneous maps on a cone $\cone$, then $\family$ is asymptotically stable if and only if $\cradius(\family)<1$.
\end{theorem}

\begin{theorem}\label{Thm_duality_for_the_cone_spectral_radius}
    Let $\family=\{f_i\}_{i\in I}$ be a bounded family of continuous, order-preserving and 1-homogeneous maps on a cone $\cone$ and let $\cradius(\family)$ be the joint cone spectral radius. Then 
    \begin{equation}
        \cradius(\family)=\inf_{\Theta\in \DMPN(\cone)}\sup_{i\in I}\Theta(f_i).
    \end{equation}
    %
    If in addition there exists some $L>0$ such that any $f_i\in \family$ is $L$-Lipschitz then
    %
    \begin{equation}
        \cradius(\family)=\inf_{\Theta\in \MPN(\cone)}\sup_{i\in I}\Theta(f_i).
    \end{equation}
    %
    In particular, if $\Sigma(\family')$ is bounded, then $\family$ admits an extremal, possibly discontinuous, monotone prenorm, where $\family'=\{f_i/\cradius(\family)\}$.   
\end{theorem}

    \begin{proof}
        If $\Theta\in \DMPN(\cone)$, then we trivially observe that since $\Theta$ is non-degenerate 
        \begin{equation}
            \cradius(\family)=\limsup_{k\rightarrow \infty}\sup_{f\in\Sigma_k(\family)}(\Theta(f))^{\frac{1}{k}}.
        \end{equation}
        Moreover, for any $f\in\Sigma_k(\family)$, $f=f_{i_1}\cdots f_{i_k}$
        \begin{equation}
            \Theta(f)\leq \Theta(f_{i_1})\cdots \Theta(f_{i_k}),
        \end{equation}
        yielding the inequality
        %
        \begin{equation}
            \cradius(\family)\leq \sup_{i\in I}\Theta(f_i).
        \end{equation}
       %
        To prove the opposite inequality, let $\psi\in \interior(\cone^*)$, then for any $\epsilon>0$ let $\family_\epsilon=\{f_i/(\cradius(\family)+\epsilon)\}_{i\in I}$ (note that since $\family$ is bounded, $\cradius(\family)$ is finite and $\family_\epsilon$ is non-zero) and define:
       %
        \begin{equation}
            \Theta_{\epsilon}(x)=\sup_{f\in \Sigma(\family_{\epsilon})}\psi(f(x)).
        \end{equation}
       %
        It is trivial to observe that  $\Theta_{\epsilon}$ is $1$-homogeneous, moreover, since $\Sigma_0(\family_\epsilon)=\{Id\}$, then $\Theta_{\epsilon}(x)\geq \psi(x)\geq 0$, with equality if and only if $x=0$. We have proved that $\Theta_\epsilon$ is a possibly discontinuous prenorm, moreover
        since any $f\in \Sigma(\family_\epsilon)$ is order preserving and $\psi(x)\geq \psi(y)$ whenever $x\cgeq y$, then $\Theta_\epsilon$ is monotone. We miss to prove that $\Theta_{\epsilon}$ is non-degenerate, first observe that:
        \begin{equation}
            \sup_x \frac{\|x\|}{\Theta_{\epsilon}(x)}\leq \sup_x \frac{\|x\|}{\psi(x)}=\frac{1}{c} <\infty
        \end{equation}
        for some $c>0$, where we have used that $\psi\in \interior(\cone^*)$.
        %
        On the other hand, by definition of the cone joint spectral radius \eqref{Cone_joint_spectral_radius}, for any $\delta>0$ there exists $N_0$ sufficiently large such that for any $k>N_0$, 
        %
        \begin{equation}        \sup_{f\in\Sigma_k(\family_\epsilon)}\|f\|^{\frac{1}{k}}\leq \frac{\cradius(\family)}{\cradius(\family)+\epsilon}+\delta.
        \end{equation}
        %
        Thus $\Sigma(\family_\epsilon)$ is bounded, i.e. there exists $\alpha>0$ such that for any $x\in\cone$, $\sup_{f\in\family_{\epsilon}}\|f(x)\|\leq \alpha\|x\|$. In particular, we have 
        %
        \begin{equation}
            \Theta_{\epsilon}(x)=\sup_{f\in \Sigma(\family_\epsilon)}\psi\big(f(x)\big)\leq \beta \sup_{f\in \Sigma(\family_\epsilon)}\|f(x)\|\leq \alpha \beta \|x\|,
        \end{equation}
        %
        for some $\beta>0$, which yields 
        %
        \begin{equation}
            \sup_{x\in\cone}\frac{\Theta_{\epsilon}(x)}{\|x\|}\leq \alpha\beta=:C<\infty,
        \end{equation}
        concluding the proof of the non-degeneracy.
        %
        %
        To conclude the theorem, observe that $
            \Theta_{\epsilon}(f(x))\leq \Theta_{\epsilon}(x)$, for any $f\in \family_{\epsilon}$ and $x\in\cone$, i.e.:
            \begin{equation}
                \Theta_\epsilon(f_i) \leq \cradius(\family)+\epsilon \qquad \forall i\in I.
            \end{equation}
          Observe also that, if $\Sigma(\family_0)=\Sigma(\family')$ is bounded, the same argument above yields the last claim of the theorem, i.e. the existence of an extremal prenorm.
          %     
          We miss to prove the Lipschitz case. In this case, recall from \Cref{Theorem_asymptotic_stability_via_JSR}, that since $\cradius(\family_\epsilon)<1$, the family $\family_\epsilon$ is asymptotically stable. In particular, for any $\epsilon>0$ there exists $N_\epsilon>0$ such that for any $k\geq N_\epsilon$ and any $f\in \Sigma_k(\family_\epsilon)$, $\psi(f(x))\leq \psi(x)$. As a consequence, we can write 
          \begin{equation}
            \Theta_{\epsilon}(x)=\sup_{f\in \cup_{k=1}^{N_\epsilon}\Sigma_k(\family_{\epsilon})}\psi(f(x)).
        \end{equation}
          Hover the last is the sup of a family of equi-Lipschitz maps and thus $\Theta_\epsilon$ is Lipschitz i.e. continuous, concluding the proof.          
    \end{proof}

\end{document}
